\documentclass[11pt,letterpaper,DIV=12]{scrartcl}

\usepackage[T1]{fontenc}
\usepackage[spanish,es-noshorthands]{babel}
\usepackage{lmodern}
\usepackage{graphicx}
\usepackage{booktabs}
\usepackage{tabularx}
\usepackage[hidelinks]{hyperref}
\usepackage{csquotes}
\usepackage[backend=biber,style=apa]{biblatex}

\addbibresource{references.bib}


\setcounter{secnumdepth}{0}

\begin{document}

\titlehead{\centering Instituto Tecnológico y de Estudios Superiores de Monterrey}
\subject{Inteligencia artificial avanzada para la ciencia de datos}
\title{Riesgos y contingencias}
\author{
  \begin{tabular}{l@{\hspace{2em}}l}
    Facundo Bautista Barbera & Emiliano Camacho Ponce \\
    A01066843 & A01712408 \\[1ex]
    Oswaldo Isaias Hernandez Santes & Iker Mejia Hernandez \\
    A01199004 & A01352358 \\[1ex]
    Alfredo Alejandro Soto Herrera & Jorge Manuel Oyoqui Aguilera \\
    A01711368 & A01711783
  \end{tabular}
}
\date{\today}

\maketitle


Dentro del proyecto pueden llegar a suscitarse situaciones en las que podamos tener retrasos o algunos fallos, a esto lo conocemos como riesgos del proyecto, para evitar que sucedan estas situaciones es necesario conocerlas y tenerlas presentes para tener cuidado con ello.

\noindent
\begin{tabularx}{\textwidth}{>{\raggedright\arraybackslash}X>{\raggedright\arraybackslash}p{1.8cm}>{\raggedright\arraybackslash}X>{\raggedright\arraybackslash}X}
  \toprule
  Riesgo & Nivel de riesgo (0 - 100) & Plan de Mitigación & Plan de Contingencia \\
  \midrule
  Dataset chico o desbalanceado & 85 & Definir un mínimo de vacas a las que tomarles foto, además de decidir cuántas fotos tomar por vaca. & Utilizar algún modelo pre entrenado con capas congeladas si es posible. \\
  \addlinespace
  Calidad irregular en las fotos & 80 & Definir una misma vista, distancia, altura, etc. para las diferentes vacas. Además de descartar fotos borrosas. & Usar data augmentation y volver a eliminar fotos irregulares. \\
  \addlinespace
  Data leakage por no separar por vacas & 80 & Separar train, validation y test por vaca y no solo al azar. & Rehacer la separación por vaca y re entrenar. \\
  \addlinespace
  El modelo aprende cosas que no son de la vaca en sí & 75 & Se puede recortar o segmentar a la vaca, además de poder implementar un bounding box a la imagen. & Segmentar y enmascarar el fondo de la imagen. \\
  \addlinespace
  Falta de tiempo al intentar resolver las 2 problemáticas & 70 & Acordar fechas de trabajo concretas, enfocándose primero en la determinación del BCS e ir midiendo el tiempo. & Entregar solamente BCS y dejar el concepto del problema de colas. \\
  \addlinespace
  Restricciones técnicas & 60 & Utilizar un iPhone de al menos un integrante del equipo, además de conseguir una GPU para el entrenamiento. & Si no se puede conseguir una cámara con 3D, usar RGB. \\
  \addlinespace
  Datos de producción de leche incompletos & 55 & En el proceso de ``Data Understanding'' investigar a fondo los datos en busca de algún problema mayor. & Imputar o excluir los días sin datos. \\
  \bottomrule
\end{tabularx}

\printbibliography

\end{document}
